\section*{Capítulo \ref{ch:g9:microbes-and-infection} --- Os micróbios e a infecção}

\begin{solution}{exo:g9:microbes-and-infection:1}
As \omterm{def:g9:microbes-and-infection:sorted}{bactérias} (\omterm{def:g5:first-look-at-cells:cell}{células} isoladas sem \omterm{def:g6:cell-unit-of-life:parts}{núcleo}), os \omterm{def:g6:decomposers-and-soil:decomposers}{fungos}
\omterm{def:g6:cell-unit-of-life:unicellular}{unicelulares} como a levedura e o resto da vida \omterm{def:g6:cell-unit-of-life:unicellular}{unicelular} da
lagoa --- e os \omterm{def:g9:microbes-and-infection:sorted}{vírus}, à parte por natureza: sem \omterm{def:g5:first-look-at-cells:cell}{célula}, sem
vida própria, apenas texto genético embalado que toma
emprestada a maquinaria de uma \omterm{def:g5:first-look-at-cells:cell}{célula} viva.
\end{solution}

\begin{solution}{exo:g9:microbes-and-infection:2}
Os trilhões de \omterm{def:g9:microbes-and-infection:sorted}{bactérias} inofensivas residentes na \omterm{def:g1:the-five-senses:sense}{pele} e no
intestino. Serviços: ajudar a \omterm{def:g4:journey-of-food:digestion}{digestão} e ocupar o terreno ---
expulsando os \omterm{prop:g9:microbes-and-infection:mostly}{patógenos} dos pontos de pouso.
\end{solution}

\begin{solution}{exo:g9:microbes-and-infection:3}
A transmissão (as gotículas de uma tosse); a entrada (um
corte, as vias aéreas); a multiplicação (divisões dobrando a
cada hora, ou \omterm{def:g5:first-look-at-cells:cell}{células} sequestradas copiando); a doença (a
latejada do dedo infeccionado, a garganta em carne viva do
resfriado).
\end{solution}

\begin{solution}{exo:g9:microbes-and-infection:4}
A \omterm{def:g9:microbes-and-infection:sorted}{bactéria} do corte se multiplica entre as \omterm{def:g5:first-look-at-cells:cell}{células} no tecido
quente e úmido, envenenando a vizinhança dela. O \omterm{def:g9:microbes-and-infection:sorted}{vírus} do
resfriado se multiplica dentro das \omterm{def:g5:first-look-at-cells:cell}{células} do revestimento
das vias aéreas, e cada \omterm{def:g5:first-look-at-cells:cell}{célula} sequestrada solta uma multidão
de cópias.
\end{solution}

\begin{solution}{exo:g9:microbes-and-infection:5}
A \omterm{def:g1:the-five-senses:sense}{pele}: a muralha, na entrada. As vassouras e os banhos dos
revestimentos (a limpeza das vias aéreas, as lágrimas, o
ácido do \omterm{def:g4:journey-of-food:tract}{estômago}): entrada e multiplicação inicial. A \omterm{prop:g9:microbes-and-infection:mostly}{flora
residente}: entrada --- o terreno ocupado. O sabonete, o
cozimento e o frio da \omterm{def:g1:healthy-habits:hygiene}{higiene}: transmissão e multiplicação.
\end{solution}

\begin{solution}{exo:g9:microbes-and-infection:6}
A transmissão --- o viajante da febre carregado nas mãos dos
médicos, da sala de autópsia até a sala de parto. Pasteur
mostrou os próprios viajantes: os \omterm{def:g6:producing-our-food:fermentation}{micróbios} como agentes do
azedamento, da doença e da \omterm{def:g9:microbes-and-infection:stages}{infecção} --- e o calor como o
freio deles.
\end{solution}

\begin{solution}{exo:g9:microbes-and-infection:7}
Os antibióticos matam de fome e quebram as \omterm{def:g9:microbes-and-infection:sorted}{bactérias} ---
\omterm{def:g5:first-look-at-cells:cell}{células} independentes, com maquinaria própria para envenenar.
O \omterm{def:g9:microbes-and-infection:sorted}{vírus} de um resfriado não tem maquinaria própria: ele roda
na \emph{sua}, e um remédio mirado no funcionamento dele mira
o das suas \omterm{def:g5:first-look-at-cells:cell}{células} --- e é por isso que o remédio da \omterm{def:g9:microbes-and-infection:sorted}{bactéria}
deixa o \omterm{def:g9:microbes-and-infection:sorted}{vírus} intacto.
\end{solution}

\begin{solution}{exo:g9:microbes-and-infection:8}
Lavar antes de comer: bloqueia a transmissão da mão para a
boca. Lavar depois do banheiro: bloqueia o ciclo de saída e
entrada da via intestinal. Assoar e cobrir a tosse (e o
\omterm{def:g1:my-body:joint}{cotovelo} do espirro): bloqueia a transmissão por gotículas na
fonte. (As regras dos dentes e da comida bloqueiam os pontos
de multiplicação da boca.)
\end{solution}

\begin{solution}{exo:g9:microbes-and-infection:9}
Patógeno: \omterm{def:g9:microbes-and-infection:sorted}{bactérias} de intoxicação alimentar, no frango.
Transmissão: o próprio prato. Entrada: o \omterm{def:g4:journey-of-food:tract}{tubo digestivo}.
Multiplicado por: três horas quentes de carro --- a regra das
\omterm{def:g6:exploring-our-environment:conditions}{condições} alimentando uma duplicação a cada meia hora, mais
ou menos. Bloqueio mais barato: a cadeia de frio --- uma
caixa térmica (ou comer na hora); o segundo mais barato,
cozinhar bem antes.
\end{solution}

\begin{solution}{exo:g9:microbes-and-infection:10}
Porque as mãos são mensageiras entre compartimentos: os
médicos de Semmelweis estavam limpos por qualquer padrão
cotidiano e ainda assim carregavam o assassino da enfermaria
do cadáver até a mãe. A lavagem que importa é a que acontece
\emph{na fronteira} --- do frango cru para a salada, do
banheiro para a mesa --- onde a via se ligaria sem isso.
\end{solution}

\begin{solution}{exo:g9:microbes-and-infection:11}
O antibiótico ataca a multiplicação --- quebrando as
\omterm{def:g9:microbes-and-infection:sorted}{bactérias} que se dividem. Usado de forma irregular, ele vira
um freio de seleção rodado pela metade: os sensíveis morrem,
a minoria resistente herda a enfermaria --- e cria o patógeno
que o próximo paciente vai encontrar.
\end{solution}

\begin{solution}{exo:g9:microbes-and-infection:12}
O texto: um pedacinho de letras genéticas, no mesmo alfabeto
de toda \omterm{def:g5:first-look-at-cells:cell}{célula} (\cref{rem:g9:genes-and-dna:universal}). A
impressora: a maquinaria de leitura e construção de uma
\omterm{def:g5:first-look-at-cells:cell}{célula} viva, que o texto que entrou vira para imprimir cópias
de si mesmo. A beirada da vida: sozinho, o \omterm{def:g9:microbes-and-infection:sorted}{vírus} não faz nada
da lista da \cref{prop:g1:living-or-not:signs} --- não se
alimenta, não cresce, não se divide; emprestado, ele roda uma
paródia de reprodução. Difícil de combater com remédio: não
existe maquinaria do \omterm{def:g9:microbes-and-infection:sorted}{vírus} para envenenar --- só a
impressora, e a impressora é você.
\end{solution}

\begin{solution}{pb:g9:microbes-and-infection:1}
\textbf{1.} Transmissão por gotículas --- tossidas e faladas
--- com a ajuda das superfícies compartilhadas nas raquetes
de tênis de mesa; entrada pelas vias aéreas.

\textbf{2.} A da multiplicação: as \omterm{def:g5:first-look-at-cells:cell}{células} sequestradas do
revestimento de cada caso imprimem cópias que semeiam os
casos seguintes --- a duplicação é a maquinaria emprestada do
\omterm{def:g9:microbes-and-infection:sorted}{vírus}, se acumulando.

\textbf{3.} A doença. Do estrago das \omterm{def:g5:first-look-at-cells:cell}{células} sequestradas do
revestimento, que estão falhando --- e, em parte, da batalha
da própria defesa, como o próximo capítulo detalha.

\textbf{4.} Porque os atos de toda a escola ligam o alcance
das gotículas de todas as turmas por uma hora: a transmissão
acompanha a aglomeração, e as duas misturas mais densas da
escola eram a via entre os compartimentos.

\textbf{5.} Alunos doentes em casa: transmissão removida na
fonte. Janelas: doses de gotículas diluídas --- transmissão.
Lavagem de mãos: via das superfícies cortada --- transmissão.
Limpeza dos equipamentos: via dos objetos compartilhados
cortada --- transmissão. (Todas mirando a mesma etapa, a mais
barata.)

\textbf{6.} ``Os antibióticos matam \omterm{def:g9:microbes-and-infection:sorted}{bactérias}; este agente é
um \omterm{def:g9:microbes-and-infection:sorted}{vírus} rodando nas próprias \omterm{def:g5:first-look-at-cells:cell}{células} das crianças --- o
remédio não encosta nele, e espalhá-lo assim mesmo só faria
rodar uma peneira de seleção sobre as \omterm{def:g9:microbes-and-infection:sorted}{bactérias} residentes de
todo mundo, criando resistência para o dia em que os
antibióticos forem realmente necessários.''

\textbf{7.} Agente: \omterm{def:g9:microbes-and-infection:sorted}{bactérias} de intoxicação alimentar.
Transmissão: a sobremesa. Entrada: o \omterm{def:g4:journey-of-food:tract}{tubo digestivo}.
Multiplicação: horas fora da geladeira --- calor, \omterm{def:g6:exploring-our-environment:conditions}{umidade},
comida e tempo, tudo concedido. Bloqueio: a geladeira ---
\omterm{def:g6:exploring-our-environment:conditions}{condições} recusadas; os dois surtos compartilham uma escola e
nada mais.

\textbf{8.} Um mapa de surto é a transmissão tornada visível:
o quem-sentava-onde desenha a via --- fileiras, duplas, atos
--- e a via, uma vez desenhada, mostra onde cortar. O
tratamento aleatório combate casos; o mapa combate a
propagação.

\textbf{9.} Transmissão: as gotículas do caso-índice, a
fileira e as duplas primeiro, e os atos para a escola toda.
Entrada: as vias aéreas, turma a turma. Multiplicação:
duplicações a cada dois dias pela primeira semana. Doença e
apagamento: as febres chegam ao pico, as medidas cortam as
vias, o conjunto de suscetíveis afina --- e a curva vira para
baixo.

\textbf{10.} Os bloqueios das medidas mataram as vias de fome
--- e o próprio conjunto era menor que a lista de chamada: os
alunos que carregavam a imunidade da cepa prima do inverno
passado não podiam ser acesos, e um surto morre quando restam
poucos que dá para pegar: a memória do próximo capítulo,
adiantada.

\textbf{11.} A de Semmelweis: cortar as vias dos mensageiros
--- isolamento em casa, mãos lavadas, raquetes limpas. A de
Pasteur: conhecer o viajante --- nomear o agente, recusar as
\omterm{def:g6:exploring-our-environment:conditions}{condições} dele e mirar o remédio no tipo dele, e não no
costume.

\textbf{12.} ``Os surtos são mais baratos de parar na
transmissão --- o sabonete, o ar, a distância e uma caixa
térmica fechada ganham de qualquer remédio que espera a
multiplicação.''
\end{solution}
